\FloatBarrier
\section{Сходимость и чувствительность выводов}
\label{sec:validation}
\subsection{Численная точность и независимые проверки}
Для двух ведущих боковых каналов номинального сценария переход
от $N=12$ к $N=13$ меняет порог не более чем на 0,00966\%.
Уточнение сетки флюенсов даёт максимальное изменение 0,000401\%,
повышение порядка пространственного ряда с 80 до 160 ---
относительное изменение $1{,}33\cdot10^{-10}$,
а ужесточение временного интегрирования --- 0,0000612\%.
При уточнении энергетической сетки NRMS изменения спектра
равен 0,0602\%, изменение усиления --- 0,0904\%, ширины ---
0,1508\%; классификация спектральных особенностей сохраняется.
Малость этих изменений свидетельствует о сходимости выбранных
наблюдаемых в пределах реализованной модели.

Отдельный усиленный расчёт продольных торцевого и бокового каналов
с $N=13$, $n_{\max}=160$, $\mathrm{rtol}=10^{-8}$,
$\mathrm{atol}=10^{-10}$ и более строгим ограничением шага даёт
пороги 49,2069 и 16,0737 $\uJ$. Отношение равно 3,0613,
тогда как исходное --- 3,0624. Следовательно, преимущество бокового
канала над торцевым не является следствием исходного шага интегрирования
или недостаточного пространственного порядка.
Однако это не означает физической точности порога 0,05\%:
при независимом сравнении полного прямого ядра с подогнанным
редуцированным откликом локальная относительная ошибка комплексного
возбуждения на отдельных энергиях всей полосы достигает 11,0\%
для торца и 3,58\% для боковой точки. Ошибка конкретного порога
и максимум ошибки комплексного спектра --- разные оценки.

Точное поле яркой моды дополнительно сверено с независимым решением
для равномерного возбуждения сфероида; относительное расхождение
коэффициента $B$ не превышает $2{,}4\cdot10^{-14}$.
Гранично-элементная проверка на сетках до 5120 панелей согласуется
с аналитическим решением качественно, но её экстраполяционная
неопределённость в близкой точке достигает примерно 14--21\%
для $B$. Поэтому она не используется как доказательство
процентной точности ближнего поля.

Для S1 отдельно проверены все шесть временных траекторий,
положительность матрицы плотности и знак интегральной работы.
Максимальное изменение полного спектра при изменении временного
окна равно $9{,}49\cdot10^{-7}$ в принятой нормировке,
разностного спектра --- $2{,}55\cdot10^{-5}$.
Ошибка преобразования падающего поля не превышает 0,1508\%,
ошибка энергетической сетки --- 0,1037\%.
Максимальный относительный хвост равен $2{,}38\cdot10^{-6}$;
минимальное собственное значение матрицы плотности составляет
$-3{,}34\cdot10^{-16}$, то есть соответствует округлению.
Независимая частотная проверка окончательного слабополевого спектра
на всех 8001 энергиях с тем же материалом $N=12$, выполненная
как с полным, так и с редуцированным пространственным ядром,
даёт максимальную нормированную разность мнимой части
поляризуемости 0,0432\%. Эти проверки относятся к восстановленному
спектру работы поля, а не к интерпретации его как чистого поглощения.

\subsection{Восстановление трёх первоначально незавершённых проверок}
Первичный архив не прошёл проверки уменьшения зазора до 0,3 нм,
уменьшения $c$ на 5\% и увеличения $a$ на 5\%. Причины лежали
на разных этапах вычислительного тракта.
Для $g=0{,}3$ нм относительное изменение пространственного ядра
при переходе от половинного к полному порядку 80 составляло
$3{,}36565\cdot10^{-5}$ при допуске $2\cdot10^{-5}$.
При порядке 160 оно уменьшается до $6{,}09\cdot10^{-11}$.
Поэтому этот отказ устраняется пространственным уточнением,
с сохранением материала, физических параметров и допуска.

При $c=11{,}875$ нм и при $a=4{,}725$ нм отказ возникал раньше:
подгонка яркого отклика с $N=13$ давала максимальную относительную
ошибку 5,0184 и 5,1062\%, соответственно, при пределе 5\%.
Это не расходимость уравнений движения. Повтор того же
детерминированного запуска использует прежний результат кэша и сам
по себе не исправляет локальный минимум оптимизации.
В дополнительном расчёте сохранённое пассивное решение использовано
только как начальное приближение; коэффициенты заново оптимизируются
для каждой новой формы по максимальной нормированной ошибке.
Неотрицательные силы, положительные частоты и затухания,
материальные допуски и проверки конечных наблюдаемых сохраняются.
Для продольного яркого отклика двух форм получены максимальные ошибки
около 4,88 и 4,87\% и NRMS около 2,44 и 2,43\%.
Тем самым исправляется процедура
численной аппроксимации, а не уравнения физической модели.
Для обеих изменённых форм пройдены пространственные и модальные
проверки всех пяти каналов. Максимальный NRMS конечного спектра
КТ среди этих десяти случаев равен 0,3443\% при прежнем допуске 5\%.

После указанных исправлений все три дополнительные проверки завершены
и приняты. Для уменьшенного зазора получены пороги 15,523 и
19,563 $\uJ$ для боковых продольного и радиального каналов.
При уменьшении $c$ до 11,875 нм они равны 13,883 и 19,558,
а при увеличении $a=b$ до 4,725 нм --- 13,583 и
18,854 $\uJ$. Во всех трёх случаях продольный канал
остаётся эффективнее радиального. Коэффициенты материала номинальной
формы при окончательном повторе малого зазора точно совпадают
с исходными; изменён только пространственный порядок.

\begin{table}[!htbp]\centering\small
\caption{Чувствительность порогов двух ведущих боковых каналов.
Каждый параметр изменяется отдельно, остальные сохраняются.
Флюенсы приведены в $\uJ$; при вариациях формы $g=0{,}5$ нм.
$E_q$ --- энергия перехода КТ, несущая остаётся 2,042 эВ.}
\label{tab:sensitivity}
\begin{tabular}{lrr}\toprule
Вариация & Боковой продольный & Боковой радиальный\\\midrule
Номинальные параметры & 16,075 & 20,229\\
$g=0{,}3$ нм & 15,523 & 19,563\\
$g=0{,}7$ нм & 16,659 & 20,961\\
$c=11{,}875$ нм & 13,883 & 19,558\\
$c=13{,}125$ нм & 18,156 & 21,102\\
$a=b=4{,}275$ нм & 18,740 & 22,092\\
$a=b=4{,}725$ нм & 13,583 & 18,854\\
$E_q-5$ мэВ & 15,757 & 20,532\\
$E_q+5$ мэВ & 16,756 & 20,382\\
$\gamma_1-20\%$ & 16,069 & 20,222\\
$\gamma_1+20\%$ & 16,081 & 20,236\\
$\gamma_\varphi-20\%$ & 16,124 & 20,707\\
$\gamma_\varphi+20\%$ & 16,074 & 19,953\\
\bottomrule\end{tabular}

\end{table}

Объединённая проверка всех 12 однофакторных вариаций принимает
сохранение порядка двух ведущих каналов
(табл.~\ref{tab:sensitivity}). Изменения формы оказывают наибольшее
влияние: продольный боковой порог меняется от 13,583 до
18,740 $\uJ$, то есть от $-15{,}50$ до $+16{,}58$\%
относительно номинального. Для изменения $g$ на $\pm0{,}2$ нм
изменения составляют $-3{,}44$ и $+3{,}63$\%.
Изменение энергии перехода на $\pm5$ мэВ даёт пороги
15,757 и 16,756 $\uJ$. Вариации $\gamma_1$ на
$\pm20$\% влияют значительно слабее, менее чем на 0,04\%
для ведущего канала, что согласуется с большой величиной $T_1$
относительно времени считывания. Вариация чистого дефазирования
на те же 20\% меняет продольный порог менее чем на 0,31\%,
радиальный --- до 2,36\%. Такое различие отражает зависимость
послеимпульсной эволюции от геометрии обратной связи.

\figone{validation_sensitivity.png}{.98}{Однофакторная чувствительность
после восстановления трёх первоначальных отказов. Показано максимальное
по контрольной группе абсолютное относительное изменение порога;
группа включает два ведущих канала и изолированную КТ.
Знак изменения можно установить по табл.~\ref{tab:sensitivity}.
Изменение несущей рассматривается отдельно.}{fig:sensitivity}

При оценке разделимости используется консервативный относительный
запас $\eta=5$\%, не являющийся статистическим доверительным
интервалом. На исходной сетке точки продольного бокового канала
$g=0{,}5;\ 0{,}75;\ 1$ нм имеют перекрывающиеся с лучшей точкой
интервалы $\Fth(1\pm\eta)$ и образуют область близкой эффективности.
Поэтому численно обоснован диапазон 0,5--1 нм на данной сетке,
а не единственный точно локализованный оптимальный зазор.
Контрольная точка $g=0{,}3$ нм относится к проверке чувствительности
и не включается задним числом в исходную сетку выбора.

Принятый критерий имеет ограниченную область действия: исходные
дискретные кандидаты плюс поочерёдные вариации параметров двух
выбранных каналов. Он не исключает появления другого лидера
среди всех пяти каналов при совместном изменении параметров,
не представляет статистический анализ ансамбля частиц и не
снимает физических ограничений субнанометровых зазоров.


\subsection{Перестройка несущей как отдельная задача управления}
При фиксированном переходе КТ изменение несущей на $\pm40$ мэВ
существенно меняет эффективность. Для 2,002 эВ боковой радиальный
порог равен 48,896 $\uJ$, а продольный боковой канал
не достигает уровня 0,5 на первой рабиевской ветви.
Для 2,082 эВ продольный боковой порог равен 32,522 $\uJ$,
а радиальный канал не достигает этого уровня на первой ветви.
Пороги изолированной КТ при этих несущих равны 147,108 и
142,725 $\uJ$, соответственно.

Из трёх проверенных несущих 2,042 эВ даёт наименьший абсолютный
порог. Но наибольший \emph{относительный} выигрыш продольного
бокового канала среди разрешённых значений получается при
2,082 эВ: около 4,39 против 3,84 при 2,042 эВ, поскольку
расстроенная изолированная КТ возбуждается значительно хуже.
Выбор метрики снова меняет формулировку оптимума.
Отсутствие первого порога одного из каналов не позволяет
утверждать неизменность полного ранжирования при перестройке несущей.
Эти три точки также не определяют непрерывный оптимум частоты.

\IfFileExists{figures/validation_carrier.png}{
\figone{validation_carrier.png}{.95}{Чувствительность порога к
энергии несущей при неизменном переходе КТ. Пропуски означают
отсутствие достижения уровня 0,5 на первой возрастающей рабиевской
ветви, а не нулевой или бесконечный вычисленный порог.}{fig:carrier}
}{}
